\documentclass[10pt,a4paper]{amsart}
\usepackage[margin=19mm]{geometry}
\usepackage{amsmath,amssymb,mathtools}
\usepackage[scr=rsfs,cal=euler]{mathalfa}
\usepackage{xfrac}
\usepackage[unicode,hidelinks,bookmarks=false]{hyperref}
\newcommand{\ie}{i.e.\ }
\newcommand{\cf}{cf.\ }
\newcommand{\resp}{respectively\ }
\newcommand{\Subsection}{\subsection}
\providecommand{\nobreakdash}{\nobreak\hbox{-}}
\setlength{\emergencystretch}{2em}
\multlinegap0pt
\usepackage{amssymb}
\usepackage{tikz}
\newtheorem{theorem}{Theorem}
\newtheorem{lemma}{Lemma}
\newtheorem{proposition}{Proposition}
\newcommand{\B}{\mathcal{B}}
\newcommand{\K}{\mathcal{K}}
\newcommand{\R}{\mathbb{R}}
\newcommand{\ga}{\gamma}
\newcommand{\la}{\lambda}
\title{Exact number of positive solutions and existence of sign-changing solutions with prescribed mass for NLS on bounded domains}
\author{Linjie Song, Wenming Zou}
\date{Source: arXiv:2603.16730v1 (CC BY 4.0)}
\begin{document}
\maketitle

\noindent\emph{Source and licence.} Exact number of positive solutions and existence of sign-changing solutions with prescribed mass for NLS on bounded domains. Version arXiv:2603.16730v1 (CC BY 4.0), distributed by arXiv under the Creative Commons Attribution 4.0 International licence, http://creativecommons.org/licenses/by/4.0/, as recorded in the arXiv licence metadata for this exact version and shown on the abstract page https://arxiv.org/abs/2603.16730v1. Full licence notice: \texttt{licenses/arxiv-2603.16730-CC-BY-4.0.txt}.\par\smallskip

\noindent\textit{Editorial scope.} Counted unit: the article's theorem thmsmall (source0.tex lines 196--204). The article's complete original proof of it is delivered with it (source0.tex lines 1728--1927, one \texttt{proof} environment of 200 lines, which the article places away from the statement). No mathematics is written, repaired or re-derived: the statement and the proof are the article's own lines, in the article's own notation, and no retained line is substituted or re-flowed.  Retained uncounted as the source-local context the proof needs: lines 160--162; lines 360--362; lines 604--609; lines 623--625; lines 685--697; lines 709--720; lines 1457--1459; lines 1475--1477; lines 1495--1501; lines 1505--1507; lines 1546--1548; lines 1552--1554; lines 1572--1574; lines 1606--1613; lines 1631--1633; lines 1656--1662; lines 1658--1660; lines 1704--1710.  Nothing was omitted from any retained span, so the package carries no omission record.  No retained line contains a citation, so the wrapper carries no bibliography and no bibliography entry.  No retained line contains a figure environment, an image-inclusion command or drawing code, so no image is delivered and the package declares no asset dependency.  Editorial formatting only, in addition: an AMS \texttt{amsart} preamble that restates the article's own declarations and macros used by the retained lines, namely the environments \texttt{theorem}, \texttt{lemma}, \texttt{proposition} on separate counters, together with the standard packages those lines need.  The retained environments print in this document as Theorem 1, Lemma 1, Proposition 1 in order of appearance, whereas the article prints the same items with the numbering of its own full document.  The complete native source of the article as distributed in the arXiv e-print for this exact version is bundled unchanged as \texttt{sources/196573/source0.tex}.
	\begin{align} \label{eqonbounddomain}
		-\Delta u + \la u = |u|^{p-2}u \quad \text{in } \Omega, \quad u \in H_0^1(\Omega), \quad \int_{\Omega}|u|^2dx = \mu.
	\end{align}

\begin{align} \label{eqgninequality}
	\|u\|_{L^p(\R^N)} \le C_{p,N}\|\nabla u\|_{L^2(\R^N)}^{\ga_p} \|u\|_{L^2(\R^N)}^{1-\ga_p},
\end{align}

\begin{lemma} \label{lemlanbelow}
	It holds that
	\begin{align*}
		\liminf_{n \to \infty}\la_n > -\infty.
	\end{align*}
\end{lemma}

\begin{lemma} \label{lemlanbounded}
	We have that $\{\la_n\}$ is bounded if and only if $\{u_n\}$ is a $L^\infty$-bounded sequence.
\end{lemma}

\begin{proposition} \label{propasm}
	As $n \to \infty$ we have
	\begin{align}
		& \tau_n^{\frac{2}{p-2}}\la_n^{\frac{N}{2} - \frac{2}{p-2}} \int_\Omega |u_n|^2dx \to \sum_{i = 1}^{k}\int_{\R^N} |V_i|^2dx, \label{eqntoinftyun2} \\
		& \tau_n^{\frac{p}{p-2}}\la_n^{\frac{N}{2} - \frac{p}{p-2}} \int_\Omega |u_n|^pdx \to \sum_{i = 1}^{k}\int_{\R^N} |V_i|^pdx, \\
		& \tau_n^{\frac{2}{p-2}}\la_n^{\frac{N}{2} - \frac{p}{p-2}}\int_\Omega |\nabla u_n|^2dx \to \sum_{i = 1}^{k}\int_{\R^N} |\nabla V_i|^2dx,
	\end{align}
    where $V_i$ is a bounded solution of
    \begin{align} \label{equationV}
    	-\Delta V + V = |V|^{p-2}V \quad \text{in } \R^N,
    \end{align}
    with $m(V) \leq \bar k$.
\end{proposition}

\begin{proposition} \label{proprelationship}
	Let $2 < p < 2^*$. The following hold:
	\begin{itemize}
		\item[(1)] If $p \neq p_c$, we have that $\{\la_n\}$ is bounded if and only if $\{E_{\tau_n}(u_n,\Omega)\}$ is bounded. Moreover,
		\begin{itemize}
			\item when $2 < p < p_c$, up to a subsequence, $\la_n \to +\infty$ if $\{\la_n\}$ is not bounded, and we have that $E_{\tau_n}(u_n,\Omega) \to -\infty$;
			\item when $p_c < p < 2^*$, up to a subsequence, $\la_n \to +\infty$ if $\{\la_n\}$ is not bounded, and we have that $E_{\tau_n}(u_n,\Omega) \to +\infty$.
		\end{itemize}
		\item[(2)] If $p = p_c$, the fact that $\{\la_n\}$ is bounded indicates that $\{E_{\tau_n}(u_n,\Omega)\}$ is bounded.
	\end{itemize}
	
\end{proposition}

	\begin{align} \label{eqdefMk}
		\mathbb{M}_k := \left\{\sum_{i=1}^ka_i\sqrt{\mu}\phi_i: \sum_{i=1}^ka_i^2 = 1\right\} \subset S_\mu(\Omega).
	\end{align}

\begin{lemma} \label{lemnonempty}
	Let $k \geq 2$. Then, for any $A \in \Sigma_{\mu}^{(k)}$ or $A \in \Sigma_{\rho,\mu}^{(k)}$ we have $A \cap \text{span}\{\phi_1,\cdots,\phi_{k-1}\}^{\bot} \neq \emptyset$.
\end{lemma}

\begin{theorem} \label{thmgenussub}
	Let $2 < p < p_c$. For $k \geq 2$ and $\mu > 0$, there exists $\tilde \delta = \tilde \delta(\mu,k) > 0$ such that when $0 < \delta < \tilde \delta$, there exists a sign-changing constrained critical point $u_{\mu,k}$ of $E(\cdot,\Omega)$ on $S_\mu(\Omega)$ at the level $\nu_{\mu,k,\delta}$. The function $u_{\mu,k}$ solves \eqref{eqonbounddomain} with $\la = \la_{\mu,k}$ for some $\la_{\mu,k} \in \R$ and $m(u_{\mu,k}) \le k+1$. Moreover, the following hold:
	\begin{itemize}
		\item when $\mu$ is fixed, for any sequence $\{\delta_k\}_{k \ge 2}$ such that $0 < \delta_k < \tilde \delta$, we have $\nu_{\mu,k,\delta_k} \to +\infty$ as $k \to +\infty$;
		\item when $k$ is fixed, for any sequence $\{\delta_\mu\}$ such that $0 < \delta_\mu < \tilde \delta$, we have $\nu_{\mu,k,\delta_\mu} \to 0$ as $\mu \to 0^+$ and $\nu_{\mu,k,\delta_\mu} \to -\infty$ as $\mu \to +\infty$.
	\end{itemize}
\end{theorem}

	\begin{align} \label{eq5.3}
		\frac12 \int_\Omega |\nabla u|^2dx \ge E(u,\Omega) \geq \frac12 \int_\Omega |\nabla u|^2dx - \frac1p C_{p,N}^p \mu^{\frac{p(1-\ga_p)}{2}}\left( \int_\Omega |\nabla u|^2dx\right) ^{\frac{p\ga_p}{2}}, \quad \forall u \in S_\mu(\Omega),
	\end{align}

    \begin{align} \label{eq5.7}
    	\nu_{\mu,k,\delta_k} \geq \inf_{S_\mu(\Omega) \cap \text{span}\{\phi_1,\cdots,\phi_{k-1}\}^{\bot}}\left( \frac12 \int_\Omega |\nabla u|^2dx - \frac1p C_{p,N}^p \mu^{\frac{p(1-\ga_p)}{2}}\left( \int_\Omega |\nabla u|^2dx\right) ^{\frac{p\ga_p}{2}} \right).
    \end{align}

    \begin{align} \label{eq7.8}
    	\nu_{\mu,k,\delta_\mu} \le \frac12 \sup_{u \in \mathbb{M}_k}\int_\Omega |\nabla u|^2dx \le \frac12\la_k(\Omega) \mu \to 0^+ \quad \text{as } \mu \to 0^+,
    \end{align}

\begin{align} \label{eqbarmu}
	0 < \mu < \bar\mu := \left( \frac{p_c}{2C_{p_c,N}^{p_c}}\right) ^{\frac{2}{p_c-2}}.
\end{align}

\begin{theorem} \label{thmgenuscri}
	Let $p = p_c$, $k \geq 2$ and let $\mu$ satisfy \eqref{eqbarmu}. Then,
	there exists $\tilde \delta = \tilde \delta(\mu,k) > 0$ such that when $0 < \delta < \tilde \delta$, there exists a sign-changing constrained critical point $u_{\mu,k}$ of $E(\cdot,\Omega)$ on $S_\mu(\Omega)$ at the level $\nu_{\mu,k,\delta}$. The function $u_{\mu,k}$ solves \eqref{eqonbounddomain} with $\la = \la_{\mu,k}$ for some $\la_{\mu,k} \in \R$ and $m(u_{\mu,k}) \le k+1$. Moreover, the following hold:
	\begin{itemize}
		\item when $\mu$ is fixed, for any sequence $\{\delta_k\}_{k \ge 2}$ such that $0 < \delta_k < \tilde \delta$, we have $\nu_{\mu,k,\delta_k} \to +\infty$ as $k \to +\infty$;
		\item when $k$ is fixed, for any sequence $\{\delta_\mu\}$ such that $0 < \delta_\mu < \tilde \delta$, we have $\nu_{\mu,k,\delta_\mu} \to 0$ as $\mu \to 0^+$.
	\end{itemize}
\end{theorem}

	\begin{align} \label{eq5.3two}
		\frac12 \int_\Omega |\nabla u|^2dx \ge E(u,\Omega) \geq \left( \frac12 - \frac1p C_{p,N}^p \mu^{\frac{p-2}{2}}\right)  \int_\Omega |\nabla u|^2dx.
	\end{align}

\begin{theorem} \label{thmgenus}
	Let $k \geq 2$, $\mu > 0$, $\rho \geq \la_k(\Omega)\mu$, $\tau > 0$ be fixed, and let us denote with $\K_c$ the set of critical points of $E$ constrained to $S_\mu(\Omega)$ at level $c$ contained in $\B_\rho^\mu$. If
	\begin{align} \label{eqconininte}
		\nu_{\rho,\mu,k,\delta} < \hat \nu_{\rho,\mu,k,\delta} := \inf_{A \in \Sigma_{\rho,\mu}^{(k)} \atop A \backslash \B_{\rho -\tau}^\mu \cap S^*(\delta) \neq \emptyset} \sup_{u \in A \backslash \B_{\rho -\tau}^\mu \cap S^*(\delta)}E(u,\Omega),
	\end{align}
    then, with sufficiently small $\delta > 0$ depending on $\rho$ and on $\mu$, the set $\K_{\nu_{\rho,\mu,k,\delta}} \cap \overline{S^*(\delta)} \neq \emptyset$ and it contains a critical point of Morse index less or equal to $k+1$. Moreover, in case $\rho \geq \la_{k+r}(\Omega)\mu$, \eqref{eqconininte} holds for $k, \cdots, k+r$, and $\nu = \nu_{\rho,\mu,k,\delta} = \cdots = \nu_{\rho,\mu,k+r,\delta}$ for some $r \ge 1$, then the genus of $\K_{\nu} \cap \overline{S^*(\delta)}$ is no less than $r+1$, so that $\K_{\nu} \cap \overline{S^*(\delta)}$ contains infinitely many critical points.
\end{theorem}

	\begin{align} 
		\nu_{\rho,\mu,k,\delta} < \hat \nu_{\rho,\mu,k,\delta} := \inf_{A \in \Sigma_{\rho,\mu}^{(k)} \atop A \backslash \B_{\rho -\tau}^\mu \cap S^*(\delta) \neq \emptyset} \sup_{u \in A \backslash \B_{\rho -\tau}^\mu \cap S^*(\delta)}E(u,\Omega),
	\end{align}

\begin{lemma} \label{lemsufficientcon}
	Let $k \geq 2$, $\mu > 0$, $\rho > 0$ satisfy
	\begin{align} \label{eqsufficientcon}
		\la_k(\Omega)\mu + \frac2p C_{p,N}^p \mu^{\frac{p(1-\ga_p)}{2}}\rho ^{\frac{p\ga_p}{2}} < \rho + \frac2p |\Omega|^{\frac{2-p}{2}}\mu^{\frac p2},
	\end{align}
	then \eqref{eqconininte} holds true for sufficiently small $\tau > 0$.
\end{lemma}

\begin{theorem} \label{thmsmall}
	Let $\Omega \subset \R^N$ be a bounded $C^1$ domain.
	\begin{itemize}
		\item[(i)] When $2 < p < p_c$, for any $\mu > 0$, \eqref{eqonbounddomain} has infinitely many sign-changing solutions $u_{\mu,1}$, $u_{\mu,2}$, $\cdots$, $u_{\mu,j}$, $\cdots$, with $\la = \la_{\mu,1}$, $\la_{\mu,2}$, $\cdots$, $\la_{\mu,j}$, $\cdots$.
		For fixed $\mu > 0$, $E(u_{\mu,j},\Omega) \to +\infty$ as $j \to +\infty$. For fixed $j$, $E(u_{\mu,j},\Omega) \to 0$ as $\mu \to 0^+$ and $E(u_{\mu,j},\Omega) \to -\infty$ as $\mu \to +\infty$. Moreover, the parameter $\displaystyle \la_{\mu,j} \to +\infty$ as $\mu \to +\infty$; the value $\displaystyle -\lim_{\mu \to 0^+}\la_{\mu,j}$ is an eigenvalue of $-\Delta$ on $\Omega$ with Dirichlet boundary condition, and particularly, $\displaystyle -\lim_{\mu \to 0^+}\la_{\mu,1} = \la_2(\Omega)$, the second Dirichlet eigenvalue.
		\item[(ii)] When $p = p_c$, there exists $\bar \mu > 0$ such that for any $0 < \mu < \bar \mu$, \eqref{eqonbounddomain} has infinitely many sign-changing solutions $u_{\mu,1}$, $u_{\mu,2}$, $\cdots$, $u_{\mu,j}$, $\cdots$, with $\la = \la_{\mu,1}$, $\la_{\mu,2}$, $\cdots$, $\la_{\mu,j}$, $\cdots$. For fixed $\mu > 0$, $E(u_{\mu,j},\Omega) \to +\infty$ as $j \to +\infty$. For fixed $j$, $E(u_{\mu,j},\Omega) \to 0$ as $\mu \to 0^+$. Moreover, the value $-\lim_{\mu \to 0^+}\la_{\mu,j}$ is an eigenvalue of $-\Delta$ on $\Omega$ with Dirichlet boundary condition, and particularly, $\displaystyle -\lim_{\mu \to 0^+}\la_{\mu,1} = \la_2(\Omega)$, the second Dirichlet eigenvalue.
		\item[(iii)] When $p_c < p < 2^*$, given a positive integer $j$, there exists $\mu^*_{p,j} > 0$ such that for any $0 < \mu < \mu^*_{p,j}$, \eqref{eqonbounddomain} has $j$ different sign-changing solutions $u_{\mu,1}$, $u_{\mu,2}$, $\cdots$, $u_{\mu,j}$, with $\la = \la_{\mu,1}$, $\la_{\mu,2}$, $\cdots$, $\la_{\mu,j}$. For fixed $i \in \{1,2,\cdots,j\}$, $E(u_{\mu,i},\Omega) \to 0$ as $\mu \to 0^+$. Moreover, the value $\displaystyle -\lim_{\mu \to 0^+}\la_{\mu,i}$ is an eigenvalue of $-\Delta$ on $\Omega$ with Dirichlet boundary condition, and particularly, $\displaystyle -\lim_{\mu \to 0^+}\la_{\mu,1} = \la_2(\Omega)$, the second Dirichlet eigenvalue.
	\end{itemize}
\end{theorem}

\begin{proof}[Proof of Theorem \ref{thmsmall}]
	We consider subcritical, critical, supercritical cases respectively.
	
	\vskip0.1in
	\emph{(1) Subcritical case with $2 < p < p_c$: }
	
	By Theorem \ref{thmgenussub}, for $k \geq 2$ and $\mu > 0$, there exists $\tilde \delta = \tilde \delta(\mu,k) > 0$ such that when $0 < \delta < \tilde \delta$, there exists a sign-changing constrained critical point $u_{\mu,k}$ of $E(\cdot,\Omega)$ on $S_\mu(\Omega)$ at the level $\nu_{\mu,k,\delta}$. Choosing $\delta_k \in (0,\tilde \delta)$, we have $\nu_{\mu,k,\delta_k} \to +\infty$ as $k \to +\infty$. Thus we can find a sequence $\{k_j\}$ with $k_1 = 2$ such that $\nu_{\mu,k_1,\delta_{k_1}} < \cdots < \nu_{\mu,k_j,\delta_{k_j}} < \cdots$. We redeonte $u_{\mu,j} = u_{\mu,k_j}$ and proved that \eqref{eqonbounddomain} has infinitely many sign-changing solutions $u_{\mu,1}, \cdots, u_{\mu,j}, \cdots$ . For fixed $\mu > 0$, we have
	\begin{align*}
		E(u_{\mu,j},\Omega) = \nu_{\mu,k_j,\delta_{k_j}} \to +\infty \quad \text{as } j \to +\infty.
	\end{align*}
	Moreover, for fixed $j$, the choice of $\delta_{k_j}$ depends on $\mu$ and we denote it by $\delta_\mu$. Then, by Theorem \ref{thmgenussub},
	\begin{align*}
		E(u_{\mu,j},\Omega) = \nu_{\mu,k_j,\delta_{\mu}} \to 0 \quad \text{as } \mu \to 0^+,
	\end{align*}
    and
    \begin{align*}
    	E(u_{\mu,j},\Omega) = \nu_{\mu,k_j,\delta_{\mu}} \to -\infty \quad \text{as } \mu \to +\infty.
    \end{align*}
    Let $\la_{\mu,j}$ be the Lagrange multiplier corresponding to $u_{\mu,j}$. As a direct consequence of Proposition \ref{proprelationship} we have $\la_{\mu,j} \to +\infty$ as $\mu \to +\infty$.

    In the following we study the limit behavior of $\la_{\mu,j}$ when $\mu \to 0^+$.

For fixed $j$, on the one hand, by \eqref{eq7.8} we have
\begin{align*}
	E(u_{\mu,j},\Omega) = \nu_{\mu,k_j,\delta_\mu} \le \frac12\la_{k_j}(\Omega) \mu.
\end{align*}
On the other hand, using \eqref{eq5.7} we get that
\begin{align*}
	E(u_{\mu,j},\Omega) = \nu_{\mu,k_j,\delta_\mu} \geq \inf_{S_\mu(\Omega) \cap \text{span}\{\phi_1,\cdots,\phi_{k_j-1}\}^{\bot}}\left( \frac12 \int_\Omega |\nabla u|^2dx - \frac1p C_{p,N}^p \mu^{\frac{p(1-\ga_p)}{2}}\left( \int_\Omega |\nabla u|^2dx\right) ^{\frac{p\ga_p}{2}} \right).
\end{align*}
Let
\begin{align*}
	f(t) = \frac12 t - \frac1p C_{p,N}^p \mu^{\frac{p(1-\ga_p)}{2}}t ^{\frac{p\ga_p}{2}}, \quad t > 0.
\end{align*}
Computing directly we get that
\begin{align*}
	f'(t) = \frac12 - \frac{\ga_p}2 C_{p,N}^p \mu^{\frac{p(1-\ga_p)}{2}}t ^{\frac{p\ga_p}{2}-1},
\end{align*}
and since $p\ga_p < 2$,
\begin{align*}
	\min_{t > 0}f(t) = f(t_\mu) = \left( \frac12 - \frac{1}{p\ga_p}\right) t_\mu \to 0 \quad \text{as } \mu \to 0^+,
\end{align*}
where
\begin{align*}
	t_\mu = \left( \frac{1}{\ga_pC_{p,N}^p}\right) ^{\frac{2}{p\ga_p-2}}\mu^{\frac{p(1-\ga_p)}{2-p\ga_p}} \to 0 \quad \text{as } \mu \to 0^+.
\end{align*}
Note that $\int_\Omega |\nabla u|^2dx \geq \la_{k_j}(\Omega)\mu$ for $u \in S_\mu(\Omega) \cap \text{span}\{\phi_1,\cdots,\phi_{k_j-1}\}^{\bot}$. By the fact
\begin{align*}
	\frac{p(1-\ga_p)}{2-p\ga_p} > 1,
\end{align*}
as $\mu \to 0^+$, it is not restrictive to assume that $\la_{k_j}(\Omega)\mu > t_\mu$. Thus
\begin{align*}
	E(u_{\mu,j},\Omega) \geq \inf_{u \in S_\mu(\Omega) \cap \text{span}\{\phi_1,\cdots,\phi_{k_j-1}\}^{\bot}}f(\int_\Omega |\nabla u|^2dx) \ge f(\la_{k_j}(\Omega)\mu).
\end{align*}
Recalling that $E(u_{\mu,j},\Omega) \le \frac12\la_{k_j}(\Omega) \mu$, we get
\begin{align*}
	\frac12\la_{k_j}(\Omega) - \frac1p C_{p,N}^p \la_{k_j}(\Omega)^{\frac{p\ga_p}{2}} \mu^{\frac{p}{2}-1} \le \mu^{-1}E(u_{\mu,j},\Omega) \le \frac12\la_{k_j}(\Omega),
\end{align*}
implying that
\begin{align*}
	\mu^{-1}E(u_{\mu,j},\Omega) \to \frac12\la_{k_j}(\Omega) \quad \text{as } \mu \to 0^+.
\end{align*}
By Lemma \ref{lemlanbounded} and Proposition \ref{proprelationship}, we have that $u_{\mu,j}$ is bounded in $L^\infty$ uniformly with respect to $\mu \to 0^+$, and so
\begin{align*}
	\int_\Omega |u_{\mu,j}|^pdx \le \|u_{\mu,j}\|_{L^\infty}^{p-2}\mu \to 0 \quad \text{as } \mu \to 0^+,
\end{align*}
thus
\begin{align*}
	\int_\Omega |\nabla u_{\mu,j}|^2dx = 2E(u_{\mu,j},\Omega) + \frac2p\int_\Omega |u_{\mu,j}|^pdx \to 0 \quad \text{as } \mu \to 0^+.
\end{align*}
By Sobolev inequality we see that
\begin{align*}
	\frac{\int_\Omega |u_{\mu,j}|^pdx}{\int_\Omega |\nabla u_{\mu,j}|^2dx} \to 0 \quad \text{as } \mu \to 0^+.
\end{align*}
Then we can conclude
\begin{align*}
	\mu^{-1}\int_\Omega|\nabla u_{\mu,j}|^2dx \to \la_{k_j}(\Omega) \quad \text{as } \mu \to 0^+.
\end{align*}
From
\begin{align*}
	\int_\Omega |\nabla u_{\mu,j}|^2 dx + \la_{\mu,j} \int_\Omega |u_{\mu,j}|^2dx = \int_\Omega |u_{\mu,j}|^p dx,
\end{align*}
it follows that
\begin{align*}
	\la_{\mu,j} = \mu^{-1}\int_\Omega |u_{\mu,j}|^p dx - \mu^{-1}\int_\Omega |\nabla u_{\mu,j}|^2 dx \to - \la_{k_j}(\Omega) \quad \text{as } \mu \to 0^+.
\end{align*}
Particularly, since $k_1 = 2$ we have $\la_{\mu,2} \to -\la_2(\Omega)$ as $\mu \to 0^+$. We complete the proof of this case.

\vskip0.1in
\emph{(2) Critical case with $p = p_c$: }

By Theorem \ref{thmgenuscri}, for $k \geq 2$ and $\mu \in (0,\bar\mu)$ where $\bar\mu$ is given in \eqref{eqbarmu}, there exists $\tilde \delta = \tilde \delta(\mu,k) > 0$ such that when $0 < \delta < \tilde \delta$, there exists a sign-changing constrained critical point $u_{\mu,k}$ of $E(\cdot,\Omega)$ on $S_\mu(\Omega)$ at the level $\nu_{\mu,k,\delta}$. Choosing $\delta_k \in (0,\tilde \delta)$, we have $\nu_{\mu,k,\delta_k} \to +\infty$ as $k \to +\infty$. Thus we can find a sequence $\{k_j\}$ with $k_1 = 2$ such that $\nu_{\mu,k_1,\delta_{k_1}} < \cdots < \nu_{\mu,k_j,\delta_{k_j}} < \cdots$. We redeonte $u_{\mu,j} = u_{\mu,k_j}$ and proved that \eqref{eqonbounddomain} has infinitely many sign-changing solutions $u_{\mu,1}, \cdots, u_{\mu,j}, \cdots$ . For fixed $\mu > 0$, we have
\begin{align*}
	E(u_{\mu,j},\Omega) = \nu_{\mu,k_j,\delta_{k_j}} \to +\infty \quad \text{as } j \to +\infty.
\end{align*}
Moreover, for fixed $j$, the choice of $\delta_{k_j}$ depends on $\mu$ and we denote it by $\delta_\mu$. Then, by Theorem \ref{thmgenuscri},
\begin{align*}
	E(u_{\mu,j},\Omega) = \nu_{\mu,k_j,\delta_{\mu}} \to 0 \quad \text{as } \mu \to 0^+.
\end{align*}

For fixed $j$, similar to the subcritical case, we have
\begin{align*}
	\left( \frac12 - \frac1p C_{p,N}^p\mu^{\frac{p}{2}-1}\right)  \la_{k_j}(\Omega) \le \mu^{-1}E(u_{\mu,j},\Omega)  \le \frac12\la_{k_j}(\Omega).
\end{align*}
Then, by \eqref{eq5.3two} it follows
\begin{align*}
	\mu^{-1}\int_\Omega |\nabla u_{\mu,j}|^2dx \to \la_{k_j}(\Omega) \quad \text{as } \mu \to 0^+.
\end{align*}
Using \eqref{eqgninequality} we obtain
\begin{align*}
	\mu^{-1}\int_\Omega|u_{\mu,j}|^pdx \le C_{p,N}^p \mu^{\frac{p}{2}-1} \mu^{-1}\int_\Omega |\nabla u_{\mu,j}|^2dx \to 0 \quad \text{as } \mu \to 0^+.
\end{align*}
Then, along the line in the subscritical case we can complete the proof.

\vskip0.1in
\emph{(3) Supercritical case with $p_c < p < 2^*$:}

Given a positive integer $j$, there exists $\mu^*_{p,j}$ such that for all $0 < \mu < \mu^*_{p,j}$, the set
\begin{align} \label{eqsetrho}
	\biggl\{\rho: \rho > \la_{j+1}(\Omega)\mu, \ \la_{j+1}(\Omega)\mu + \frac2p C_{p,N}^p \mu^{\frac{p(1-\ga_p)}{2}}\rho ^{\frac{p\ga_p}{2}} < \rho + \frac2p |\Omega|^{\frac{2-p}{2}}\mu^{\frac p2}\biggr\}
\end{align}
is not empty. In fact, there exists $\mu^*_{p,j}$ such that for any $0 < \mu < \mu^*_{p,j}$, we can take $s = s(\mu) > 0$ such that
\begin{align*}
	\frac2p C_{p,N}^p \mu^{\frac{p}{2}}(\la_{j+1}(\Omega) + s) ^{\frac{p\ga_p}{2}} < s\mu + \frac2p |\Omega|^{\frac{2-p}{2}}\mu^{\frac p2}
\end{align*}
and thus $\rho = (\la_{j+1}(\Omega) + s)\mu$ belongs to the set defined in \eqref{eqsetrho}, showing that this set is not empty. For $\mu \in (0,\mu^*_{p,j})$ we take $\rho = \rho_\mu$ in the set defined in \eqref{eqsetrho}. By Lemma \ref{lemsufficientcon}, there exists $\tau > 0$ such that \eqref{eqconininte} holds true for $k = 2,\cdots,j+1$. By Theorem \ref{thmgenus}, for any $k \in \{2,\cdots,j+1\}$, there exists a sign-changing constrained critical point $\tilde u_{\mu,k}$ of $E(\cdot,\Omega)$ on $S_\mu(\Omega)$ at the level $\nu_{\rho,\mu,k,\delta}$, whose Morse index is less or equal to $k+1$.

We firstly assume that $\nu_{\rho,\mu,2,\delta} < \nu_{\rho,\mu,3,\delta} < \cdots < \nu_{\rho,\mu,j+1,\delta}$. Denote $u_{\mu,1} = \tilde u_{\mu,2}, \cdots, u_{\mu,j} = \tilde u_{\mu,j+1}$, and these are $j$ different sign-changing solutions of \eqref{eqonbounddomain} with $\la = \la_{\mu,1}$, $\la_{\mu,2}$, $\cdots$, $\la_{\mu,j}$. For fixed $i \in \{1,\cdots,j\}$,  denoting $\rho = \rho_\mu$ and $\delta = \delta_\mu$, by the fact that $\mathbb{M}_{i+1}  \in \Sigma_{\rho_\mu,\mu}^{(i+1)}$, it follows that
\begin{align} \label{eq7.18}
	E(u_{\mu,i},\Omega) = \nu_{\rho_\mu,\mu,i+1,\delta_\mu} \le \frac12 \sup_{u \in \mathbb{M}_{i+1}}\int_\Omega |\nabla u|^2dx \le \frac12\la_{i+1}(\Omega) \mu \to 0^+ \quad \text{as } \mu \to 0^+,
\end{align}
where $\mathbb{M}_{i+1}$ is defined in \eqref{eqdefMk}. Thus we get $\limsup_{\mu \to 0^+}\nu_{\rho_\mu,\mu,i+1,\delta_\mu} \le 0$. In the following we prove that $\liminf_{\mu \to 0^+}\nu_{\rho_\mu,\mu,i+1,\delta_\mu} \ge 0$. By negation, we suppose that there exists $\{\mu_n\}$ such that $\mu_n \to 0$ and $\lim_{n \to \infty}\nu_{\rho_n,\mu_n,i+1,\delta_n} < 0$, where $\rho_n = \rho_{\mu_n}$ and $\delta_n = \delta_{\mu_n}$. Let us take a sequence $\{u_n\} = \{u_{\mu_n,i}\} \subset S_{\mu_n}(\Omega)$ such that $E(u_n,\Omega) = \nu_{\rho_n,\mu_n,i+1,\delta_n}$, $u_n$ solves \eqref{eqonbounddomain} with $\la = \la_n$, $\mu = \mu_n \to 0$, and $m(u_n) \le i+2$. Similar to the proof of Theorem \ref{thmgenussub}, by Lemma \ref{lemlanbelow} and Lemma \ref{lemlanbounded} we get that $\la_n \to +\infty$. Following the arguments in proving Proposition \ref{proprelationship}, by using Proposition \ref{propasm} we can prove that $E(u_n,\Omega) \to +\infty$, contradicting \eqref{eq7.18}. This shows that $E(u_{\mu,i},\Omega) \to 0$ as $\mu \to 0^+$. Moreover, by Lemma \ref{lemnonempty} and \eqref{eq5.3} it holds that
\begin{align} \label{eqnuineq}
	\nu_{\rho_\mu,\mu,i+1,\delta_\mu} \geq \inf_{\B_{\rho_\mu}^\mu \cap \text{span}\{\phi_1,\cdots,\phi_{i}\}^{\bot}}\left( \frac12 \int_\Omega |\nabla u|^2dx - \frac1p C_{p,N}^p \mu^{\frac{p(1-\ga_p)}{2}}\left( \int_\Omega |\nabla u|^2dx\right) ^{\frac{p\ga_p}{2}} \right).
\end{align}
Recall that $f$ is introduced in the subcritical case by
\begin{align*}
	f(t) = \frac12 t - \frac1p C_{p,N}^p \mu^{\frac{p(1-\ga_p)}{2}}t ^{\frac{p\ga_p}{2}}, \quad t > 0,
\end{align*}
and
\begin{align*}
	f'(t) = \frac12 - \frac{\ga_p}2 C_{p,N}^p \mu^{\frac{p(1-\ga_p)}{2}}t ^{\frac{p\ga_p}{2}-1}.
\end{align*}
In this case, $p\ga_p > 2$, thus $f(t)$ is strictly decreasing in $t \ge t_\mu$,
where
\begin{align*}
	t_\mu = \left( \frac{1}{\ga_pC_{p,N}^p}\right) ^{\frac{2}{p\ga_p-2}}\mu^{\frac{p(1-\ga_p)}{2-p\ga_p}} \to 0 \quad \text{as } \mu \to 0^+.
\end{align*}
As discussed at the start of this case, $\rho_\mu$ can be chosen as $\rho_\mu = (\la_{i+1}(\Omega)+s(\mu))\mu$ where
\begin{align*}
	s(\mu) \to 0 \quad \text{and} \quad s(\mu)\mu^{\frac{2-p}{2}} \to +\infty \quad \text{as } \mu \to 0^+.
\end{align*}
Note that $\la_{i+1}(\Omega)\mu \le \int_\Omega |\nabla u|^2dx \le (\la_{i+1}(\Omega)+s(\mu))\mu$ for $u \in \B_{\rho_\mu}^\mu \cap \text{span}\{\phi_1,\cdots,\phi_{i}\}^{\bot}$. By the fact
\begin{align*}
	\frac{p(1-\ga_p)}{2-p\ga_p} > 1,
\end{align*}
as $\mu \to 0^+$, it is not restrictive to assume that $\la_{i+1}(\Omega)\mu > t_\mu$. Thus \eqref{eqnuineq} yields that
\begin{align*}
	\nu_{\rho_\mu,\mu,i+1,\delta_\mu} \geq f((\la_{i+1}(\Omega)+s(\mu))\mu) = \frac12\la_{i+1}(\Omega)\mu + o(\mu),
\end{align*}
where $\mu^{-1}o(\mu) \to 0$ as $\mu \to 0^+$. Together with \eqref{eq7.18}, we conclude that
\begin{align*}
	\mu^{-1}E(u_{\mu,i},\Omega) \to \frac12\la_{i+1}(\Omega) \quad \text{as } \mu \to 0^+.
\end{align*}
Then, following the discussions in the proof of subcritical case, we can obtain that $\la_{\mu,i} \to -\la_{i+1}(\Omega)$, particularly, $\la_{\mu,1} \to -\la_2(\Omega)$ as $\mu \to 0^+$.

Next, we assume that $\nu_{\rho,\mu,k,\delta} = \nu_{\rho,\mu,k+1,\delta}$ for some $k \in \{2,\cdots,j\}$. By Theorem \ref{thmgenus}, \eqref{eqonbounddomain} has $j$ different sign-changing solutions $u_{\mu,1}$, $u_{\mu,2}$, $\cdots$, $u_{\mu,j}$, with $\la = \la_{\mu,1}$, $\la_{\mu,2}$, $\cdots$, $\la_{\mu,j}$, where $u_{\mu,1}$ is at level $\nu_{\rho,\mu,2,\delta}$, $u_{\mu,2}, \cdots, u_{\mu,j}$ are at level $\nu = \nu_{\rho,\mu,k,\delta} = \nu_{\rho,\mu,k+1,\delta}$, and $m(u_{\mu,1}) \le 3$, $m(u_{\mu,2}) \le k+1$. As proved above, we have
\begin{align*}
	\mu^{-1}\nu_{\rho_\mu,\mu,2,\delta_\mu} \to \frac12\la_2(\Omega) \quad \text{and} \quad \mu^{-1}\nu \to \frac12\la_{k}(\Omega) \quad \text{as } \mu \to 0^+,
\end{align*}
which implies that
\begin{align*}
	E(u_{\mu,i},\Omega) \to 0 \quad \text{as } \mu \to 0^+
\end{align*}
for any $i \in \{1,2,\cdots,j\}$. Moreover, we have that $\la_{\mu,1} \to -\la_2(\Omega)$ and $\la_{\mu,2} \to -\la_k(\Omega)$ as $\mu \to 0^+$. In the following we consider $i \in \{3,\cdots,j\}$. By the boundedness of $u_{\mu,i}$ in $H_0^1(\Omega)$ uniformly with respect to $\mu \to 0^+$ and by Gagliardo-Nirenberg inequality \eqref{eqgninequality} we get
\begin{align*}
	\int_\Omega|u_{\mu,i}|^pdx \to 0 \quad \text{as } \mu \to 0^+,
\end{align*}
and so
\begin{align*}
	\int_\Omega |\nabla u_{\mu,i}|^2dx = 2E(u_{\mu,i},\Omega) + \frac2p\int_\Omega |u_{\mu,i}|^pdx \to 0 \quad \text{as } \mu \to 0^+.
\end{align*}
By Sobolev inequality we see that
\begin{align*}
	\frac{\int_\Omega |u_{\mu,i}|^pdx}{\int_\Omega |\nabla u_{\mu,i}|^2dx} \to 0 \quad \text{as } \mu \to 0^+.
\end{align*}
Thus we can conclude
\begin{align*}
	\mu^{-1}\int_\Omega|\nabla u_{\mu,i}|^2dx \to \la_k(\Omega) \quad \text{as } \mu \to 0^+.
\end{align*}
From
\begin{align*}
	\int_\Omega |\nabla u_{\mu,i}|^2 dx + \la_{\mu,i} \int_\Omega |u_{\mu,i}|^2dx = \int_\Omega |u_{\mu,i}|^p dx,
\end{align*}
it follows that
\begin{align*}
	\la_{\mu,i} = \mu^{-1}\int_\Omega |u_{\mu,i}|^p dx - \mu^{-1}\int_\Omega |\nabla u_{\mu,i}|^2 dx \to -\la_k(\Omega) \quad \text{as } \mu \to 0^+.
\end{align*}
The proof is complete.
\end{proof}

\end{document}
